% O resumo é obrigatório, em português e inglês. Este comando também gera
% automaticamente a referência para o próprio documento, conforme as normas
% sugeridas da USP
\begin{resumo}{port}
\textbf{Contexto:} Projetos de software livre constituem parte significativa do mercado de desenvolvimento de software. Neste contexto, para que os projetos se mantenham ativos, contribuições voluntárias são esperadas. No entanto, novos contribuidores, ao tentarem ingressar voluntariamente em projetos de software livre, costumam enfrentar um conjunto de barreiras que acabam por dificultar o processo de contribuição. Tais barreiras tendem a levar novatos a desistirem de contribuir com projetos de software livre e, por esta razão, são amplamente estudadas pela literatura. Entre as dificuldades já constatadas, existem barreiras associadas as documentações dos projetos de software livre que, muitas vezes, não atendem as necessidades que novos contribuidores tem durante o processo de contribuição.

\textbf{Objetivo:} Almejando contribuir com a avaliação de documentações de projetos de software livre, propomos como objetivo desta pesquisa a construção de um modelo de classificação, capaz de identificar trechos em documentações de projetos de software livre que sejam relevantes a novatos durante o processo de contribuição. A ideia é que esta ferramenta, ao fim deste estudo, auxilie novos contribuidores a compreenderem aspectos essenciais dos projetos aos que eles almejem contribuir, bem como permitam uma melhor avaliação dos projetos quanto a qualidade das documentações que estes produzem para novatos. 

\textbf{Método:} A partir de um experimento a ser realizado com novatos provenientes de projetos de software livre, um modelo de classificação de trechos relevantes a novos contribuidores será construído. Utilizando como base fundamental algoritmos de aprendizado, este modelo receberá como entrada parágrafos de documentações de projetos, e deverá identificar para cada um dos parágrafos, um conjunto de categorias de informação relevantes a novos contribuidores, que serão definidas em um etapa preliminar. O modelo de classificação, bem como os dados e resultados obtidos durante a execução deste método, serão analisados sob diferentes perspectivas a partir de duas etapas, uma de avaliação estatística, e outra através de questionários, a serem aplicados com novos e experimentes contribuidores dentro do contexto de software livre. 

\textbf{Resultados esperados:} Esperamos com este estudo responder, se é possível identificar trechos de documentação relevantes a novos contribuidores em software livre utilizando algoritmos de classificação, além de apresentar quais características são mais relevantes para identificação destes trechos, e a aceitação das predições feitas pelo modelo de classificação por parte de contribuidores neste contexto.
\end{resumo}

% O resumo é obrigatório, em português e inglês. Este comando também gera
% automaticamente a referência para o próprio documento, conforme as normas
% sugeridas da USP
\begin{resumo}{eng}
\textbf{Context:} Open source projects constitute a significant part of the software development market. In this context, for projects to remain active, voluntary contributions are expected. However, new contributors, when voluntarily trying to join open source projects, often face a set of barriers that compromises the contribution process. Such barriers tend to lead newcomers to give up contributing to these projects and, for this reason, are widely studied in the literature. Among the difficulties already noted by literature, there are barriers associated with the documentation of open source projects that often do not meet the needs that new contributors have during the contribution process.

\textbf{Objective:} Aiming to contribute to the documentation of open source projects, we propose the development of a classification model, capable of identifying snippets in open source projects documentation that are relevant to newcomers during the process. contribution. The idea is that this tool, at the end of this study, will help newcomers to understand essential aspects of the projects they aim to contribute, minimizing potential difficulties faced during the contribution process.

\textbf{Method:} From an experiment to be conducted with newcomers from open source projects, and researchers in software engineering, a model of classification of categories relevant to new contributors will be built. Using machine learning concepts as a fundamental basis, this model will receive as input project documentation paragraphs, and should identify for each paragraph a set of categories of information relevant to new contributors, which will be defined at a preliminary stage. The classification model, as well as the data and results obtained during the execution of this method, will be analyzed from different perspectives in a statistical evaluation stage, and through questionnaires, to be applied with new and core contributors in the context of open source software.

\textbf{Expected Results:} We hope with this study to answer, whether it is possible to identify relevant documentation snippets to new open source contributors using classification algorithms, and to present which features are most relevant for identifying these snippets, and the acceptance of predictions by contributors in this context.
\end{resumo}
